\DocumentMetadata{
  pdfversion=2.0,
  pdfstandard=ua-2,
  lang=en-US,
  tagging=on,
  tagging-setup={math/setup=mathml-SE,tabsorder=structure}
}
\documentclass[11pt,letterpaper]{article}
\usepackage[margin=0.68in,includefoot]{geometry}
\usepackage{newcomputermodern}
\usepackage{amsmath,amscd,mathtools}
\usepackage{tikz,adjustbox}
\providecommand{\square}{\mdlgwhtsquare}
\usepackage[hidelinks]{hyperref}
\hypersetup{
  pdftitle={Algebra PhD Exam, January 2022},
  pdfauthor={University of Florida Department of Mathematics},
  pdfsubject={Graduate mathematics examination},
  pdfdisplaydoctitle=true,
  bookmarksopen=true
}
\setcounter{secnumdepth}{0}
\setlength{\parindent}{0pt}
\setlength{\parskip}{0.28em}
\setlength{\emergencystretch}{3em}
\setlength{\leftmargini}{2.15em}
\setlength{\leftmarginii}{2.65em}
\setlength{\leftmarginiii}{3em}
\pagestyle{empty}
\raggedbottom
\allowdisplaybreaks
\NewTaggingSocketPlug{block/list/label}{examproblem}{%
  \tagstructbegin{tag=Lbl}\tagstructend%
  \tagstructbegin{tag=\LBody}%
  \tagstructbegin{tag=\ProblemHeadingTag,title-o={\ProblemHeadingText}}%
  \tagmcbegin{tag=\ProblemHeadingTag,actualtext={\ProblemHeadingText}}%
  #2%
  \tagmcend%
  \tagstructend%
}
\newenvironment{examproblems}[2]{%
  \def\ProblemHeadingTag{#2}%
  \AssignTaggingSocketPlug{block/list/label}{examproblem}%
  \begin{enumerate}%
  \setcounter{enumi}{#1}%
  \renewcommand{\labelenumi}{\textbf{\theenumi.}}%
  \setlength{\topsep}{0.35em}%
  \setlength{\partopsep}{0pt}%
  \setlength{\itemsep}{0.48em}%
  \setlength{\parsep}{0.18em}%
}{\end{enumerate}}
\newenvironment{examsubparts}{%
  \AssignTaggingSocketPlug{block/list/label}{default}%
  \begin{enumerate}%
  \setlength{\topsep}{0.18em}%
  \setlength{\partopsep}{0pt}%
  \setlength{\itemsep}{0.16em}%
  \setlength{\parsep}{0.1em}%
}{\end{enumerate}}
\newenvironment{examinstructions}{%
  \par\begingroup\itshape
}{%
  \par\endgroup\vspace{0.18em}
}
\makeatletter
\renewcommand\subsection{\@startsection{subsection}{2}{0pt}%
  {-1.25ex plus -.25ex minus -.1ex}%
  {0.55ex plus .1ex}%
  {\normalfont\large\bfseries}}
\renewcommand{\@maketitle}{%
  \newpage\null\vskip -1em%
  {\footnotesize\bfseries
    \href{https://www.ufl.edu/}{University of Florida}, \href{https://math.ufl.edu/}{Department of Mathematics}\par}%
  \vskip -0.5em%
  \tagstructbegin{tag=H1,title={Algebra PhD Exam, January 2022}}%
  \tagpdfparaOff%
  \tagmcbegin{tag=H1}%
    {\LARGE\bfseries\href{https://gma.math.ufl.edu/exams/algebra-phd/}{Algebra PhD Exam}, January 2022\par}%
  \tagmcend%
  \tagpdfparaOn%
  \tagstructend%
  \vskip 0.25em%
  \tagmcbegin{artifact}%
    \hrule height 1pt%
  \tagmcend%
  \vskip 1em%
}
\makeatother
\title{Algebra PhD Exam, January 2022}
\author{}
\date{}
\begin{document}
\maketitle
\thispagestyle{empty}
\begin{examinstructions}
Answer seven problems. If more than seven problems are answered only the first seven will be graded. Write your answers clearly in complete English sentences.
\end{examinstructions}
\begin{examinstructions}
Results from lectures or textbooks may be used without proof (within reason – don’t state results equivalent to the problem), but must be clearly stated.
\end{examinstructions}
\begin{examproblems}{0}{H2}
\def\ProblemHeadingText{Problem 1.}
\item
Suppose~\(K\) is a field of characteristic different from~\(2\), \(a \in K\) is not a square and \(b \in K\) is not a square in \(K(\sqrt{a})\). Compute the degree of \(L=K(\sqrt{a},\sqrt{b})\) over~\(K\) and determine all subfields of~\(L\) containing~\(K\).
\def\ProblemHeadingText{Problem 2.}
\item
Suppose~\(A\) is an integral domain with fraction field~\(K\), \(L/K\) is a field extension and \(x \in L\). Show that~\(x\) is integral over~\(A\) if and only if there is a finitely generated~\(A\)-module \(M \subseteq L\) such that \(xM \subseteq M\).
\def\ProblemHeadingText{Problem 3.}
\item
Suppose~\(C\) is a category and~\(X\) and~\(Y\) are objects of~\(C\).
\begin{examsubparts}
\item[\textnormal{(i)}]
Define what it means for an object of~\(C\) to be a product of~\(X\) and~\(Y\).
\item[\textnormal{(ii)}]
Show that a product of~\(X\) and~\(Y\) is unique up to isomorphism.
\end{examsubparts}
\def\ProblemHeadingText{Problem 4.}
\item
Suppose~\(G\) is a finite nilpotent group. Show that if~\(H\) is a proper subgroup of~\(G\), it is a proper subgroup of its normalizer in~\(G\).
\def\ProblemHeadingText{Problem 5.}
\item
Let~\(R\) be a integral domain.
\begin{examsubparts}
\item[\textnormal{(i)}]
Define what it means for a fractional ideal of~\(R\) to be invertible.
\item[\textnormal{(ii)}]
Show that an invertible ideal is a projective~\(R\)-module.
\end{examsubparts}
\def\ProblemHeadingText{Problem 6.}
\item
Suppose~\(R\) is a ring with identity (not necessarily commutative), \(M\) is a right~\(R\)-module and
\[
N' \to N \to N'' \to 0
\]
 is an exact sequence of left~\(R\)-modules. Show that
\[
M \otimes_R N' \to M \otimes_R N \to M \otimes_R N'' \to 0
\]
 is an exact sequence of abelian groups.
\def\ProblemHeadingText{Problem 7.}
\item
Suppose~\(F\) is a free group on a set~\(S\), \(T\) is a subset of~\(S\) and \(H \subseteq F\) is the smallest normal subgroup of~\(F\) containing~\(T\). Show that \(F/H\) is a free group on the set \(S \setminus T\).
\def\ProblemHeadingText{Problem 8.}
\item
Suppose~\(R\) is a commutative ring with identity and~\(I\) is an ideal contained in the Jacobson radical of~\(R\). Prove Nakayama’s lemma: if~\(M\) is a finitely generated~\(R\)-module such that \(IM=M\) then \(M=0\).
\def\ProblemHeadingText{Problem 9.}
\item
Suppose~\(A\) is a Noetherian commutative ring with identity.
\begin{examsubparts}
\item[\textnormal{(i)}]
Show that if \(f:A \to B\) is a surjective homomorphism then~\(B\) is Noetherian.
\item[\textnormal{(ii)}]
Recall that a multiplicative system in~\(A\) is a nonempty subset closed under multiplication. Show that if \(S \subset A\) is a multiplicative system then \(S^{-1}A\) is Noetherian.
\end{examsubparts}
\def\ProblemHeadingText{Problem 10.}
\item
Suppose~\(R\) is a nontrivial commutative ring with identity and \(S \subset R\) is a multiplicative system not containing~\(0\) (see the previous problem for the definition).
\begin{examsubparts}
\item[\textnormal{(i)}]
Show that there is a prime ideal \(\mathfrak p \subset R\) such that \(\mathfrak p \cap S=\varnothing\). Hint: Apply Zorn’s lemma on the ideals \(I \subset R\) such that \(I \cap S=\varnothing\).
\item[\textnormal{(ii)}]
Use this to show that the intersection of all prime ideals of~\(R\) is the set of nilpotent elements of~\(R\).
\end{examsubparts}
\def\ProblemHeadingText{Problem 11.}
\item
This problem has two parts.
\begin{examsubparts}
\item[\textnormal{(i)}]
State the Jacobson density theorem for semisimple rings.
\item[\textnormal{(ii)}]
Use it to show that if~\(K\) is a field and~\(A\) is a simple~\(K\)-algebra of finite dimension then there is a division~\(K\)-algebra~\(D\) such that \(A \simeq M_n(D)\) for some~\(n\).
\end{examsubparts}
\end{examproblems}
\end{document}
